\subsection{The contact strip}
\label{sec:direct-bounded-contact}

We treat the contact strip, with fixed \(L<\infty\),
\begin{equation}
\label{eq:direct-contact-region}
 q_1>0,\qquad q_0=\widehat q_0 q_1^2,\qquad |\widehat q_0|\le L .
\end{equation}
The limits of \(q_1\) and \(q_2\) are independent, including
\(q_2=0\); constants may depend on \(L\) and \(C_0\).
The difficulty is to compare the two passages uniformly through
contact, without relating the sizes of \(q_1\) and \(q_2\).
The paired identity gives opposite derivative signs at the two ends
of the height interval. On the compact interval between them,
\cref{lem:direct-contact-scattering-ratio} gives a strictly decreasing
ratio of the limiting scattering derivatives. We transfer its derivative
margin to the endpoint maps by a \(C^2\) estimate. Every middle zero
of \(\mathcal L\) then crosses upwards in \(T\), and the outer signs
complete the count on \(\mathcal H_p\).

\begin{proposition}[Bounded contact offsets]
\label{prop:direct-bounded-contact}
For every fixed \(L\), the return displacement
\eqref{eq:central-displacement-new} has at most two
distinct zeros on the height interval \(\cH_p\) in
\eqref{eq:direct-contact-region}, for all sufficiently small parameters.
\end{proposition}

\subsubsection{Barrier and tail estimates}

With the contact scales of \cref{tab:direct-scale-guide}, put
\[
 d_{\rm c}=\max\{|q_2|,\sqrt{q_1}\},\qquad
 z=q_1^2 Z,\qquad h=q_1^2 \widehat h,\qquad
 |\eta_\sigma|=q_1u_\sigma,\qquad
 \mathcal U_\sigma=u_\sigma^2/2 .
\]
The rescaled equations are
\begin{equation}
\label{eq:direct-contact-equations}
 (\mathcal U_\sigma)_Z
 =B(\sigma q_1u_\sigma,q_1^2Z)
  +\frac{-\widehat q_0-\sigma u_\sigma-q_2u_\sigma^2
               -\sigma q_1u_\sigma^3R_q(\sigma q_1u_\sigma)}{Z}.
\end{equation}
They are the graphs of the analytic outward time equations
\begin{equation}
\label{eq:direct-contact-time}
 (u_\sigma)_\tau=-\widehat q_0-\sigma u_\sigma-q_2u_\sigma^2
               -\sigma q_1u_\sigma^3R_q+ZB,\qquad
 Z_\tau=u_\sigma Z .
\end{equation}
The central point is \(u_\sigma=0,\ Z=\widehat h\). Its radial flux is
\[
 f(\widehat h)=\widehat h B(0,q_1^2\widehat h)-\widehat q_0,\qquad
 f_{\widehat h}=B(0,q_1^2\widehat h)
   +q_1^2\widehat h B_z(0,q_1^2\widehat h)\ge b_0.
\]
Thus all \(\widehat h>0\) pass when \(\widehat q_0\le0\); when \(\widehat q_0>0\)
the passing interval is \(\widehat h>\widehat h_0\), where \(f(\widehat h_0)=0\).

The negative graph has the uniform upper barrier
\begin{equation}
\label{eq:direct-contact-cap}
 U_-(z)\le C(d_{\rm c}^2+z),\qquad |\eta_-(z)|\le C(d_{\rm c}+\sqrt z).
\end{equation}
Indeed, for \(x\ge Kd_{\rm c}\), the coefficient bounds give
\[
 q(-x)\ge cx^3-d_{\rm c}x^2-q_1x-Lq_1^2\ge0.
\]
After division by \(x^3\), the three adverse terms are at most
\(K^{-1},K^{-2},Ld_{\rm c}/K^3\). Choose \(K\) large and then \(d_{\rm c}\)
small. On the line \(U=C(d_{\rm c}^2+z)\), with \(\sqrt{2C}>K\),
the graph derivative \(B-q(-\sqrt{2U})/z\) is at most
\(\sup B<C\). The line starts above the central point, so first contact
proves \eqref{eq:direct-contact-cap}.

To continue from bounded rescaled central heights to the fixed section,
we first obtain the following cone bounds. For each fixed bounded set
of rescaled central heights, they hold for a sufficiently large fixed
\(Z_0\):
\begin{equation}
\label{eq:direct-contact-cones}
 cz\le U_\pm(z),\qquad U_+(z)\le Cz,\qquad
 z_0:=q_1^2Z_0\le z\le z_*.
\end{equation}
For continuation on a fixed interval \(\widehat h\le Z\le Z_1\),
first bound the positive radius by \(M_{L,Z_1}\). At a fixed
\(M>2(L+(\sup B)Z_1+1)\), the first right side of
\eqref{eq:direct-contact-time} is negative, after reducing
\(|q_2|\) and \(q_1\). Hence the graph cannot first reach \(M\).
Choose \(Z_1\) larger than twice these rescaled central heights and
so large that \(L/Z_1<b_0/4\). On the lower comparison line
\(\mathcal U=m(Z-Z_1)\), the adverse terms in
\eqref{eq:direct-contact-equations} are bounded by
\[
 L/Z_1+C\sqrt{m/Z_1}+2d_{\rm c}m+C m^{3/2}\sqrt{z_*}.
\]
Choose \(m>0\) small after \(Z_1\), and use the common-cone choice of
\(z_*\) already fixed before the regional thresholds.  Reducing
\(d_{\rm c}\) and \(q_1\) after these choices makes the graph derivative
greater than \(m\).  This proves a lower bound for both
signs. For the positive upper bound use \(W=C_2Z+C_3\), choosing
\(C_3\) to dominate the energy at \(Z_1\) and to have
\(\sqrt{2W}\ge2L\). The term \(-\widehat q_0-\sqrt{2W}\) is then
nonpositive, and the remaining derivative is at most
\[
 \sup B+2d_{\rm c}W/Z+Cq_1W^{3/2}/Z<C_2
\]
for \(C_2\) fixed large and the same successive reductions.
Enlarge \(Z_0\) relative to \(Z_1,C_3\). These comparisons give
\eqref{eq:direct-contact-cones}, with fixed cone constants, up
to \(z_*\). Together with \eqref{eq:direct-contact-cap}
they keep both graphs inside the box in the coordinates \((\eta,z)\).

The negative graph may remain above the upper cone for part of the
passage. To estimate the resulting loss in the paired identity, set
\(v_0=|\eta_-(z_0)|/q_1\) and
\(t_h=\log(Z_0/\widehat h)\). For \(z\ge z_0\), discard the favorable
negative cubic term and use
\(q_1\sqrt{2U}\le U/8+4q_1^2\). Taking \(Z_0\ge L+4\)
and then \(d_{\rm c}\) small yields
\begin{equation}
\label{eq:direct-contact-upper-comparison}
 U_-'\le C+\tfrac12 U_-/z,\qquad
 U_-(z)\le Cz+U_-(z_0)(z/z_0)^{1/2}.
\end{equation}
The second term is at most \(z\) once
\(z\ge z_0\max\{1,(U_-(z_0)/z_0)^2\}\).
Below that height use \eqref{eq:direct-contact-cap}; above it use
the resulting upper cone. Splitting the integral, also when this
height exceeds \(z_*\), gives
\begin{equation}
\label{eq:direct-contact-tail-loss}
 \int_{z_0}^{z_*}\frac{|\eta_-(z)|}{z}\,dz
       \le C\sqrt{z_*}+Cd_{\rm c}\log(1+v_0).
\end{equation}

The tail loss must be smaller than the gain in the paired identity
at small central heights. We therefore express its last logarithm in
terms of the central height. In the
negative equation use \(t=\log(Z/\widehat h)\) and
\(\mathcal V=\sqrt{1+u_-^2}\). Dropping the favorable cubic
term, and using \(|\widehat q_0|\le L\), gives
\[
 \mathcal V_t\le C_L+d_{\rm c}\mathcal V+(\sup B)\widehat h \e^t,\qquad
 \mathcal V(t)\le \e^{d_{\rm c}t}(1+C_Lt)+C\widehat h \e^t .
\]
In particular,
\begin{equation}
\label{eq:direct-contact-source-growth}
 \log(1+v_0)\le C_{L,Z_0}+d_{\rm c}t_h+\log(1+t_h).
\end{equation}
This bounds the tail loss in terms of \(t_h\), uniformly
in \(d_{\rm c}|\log q_1|\).

\subsubsection{The two ends of the height interval}

For the fixed contact width \(L\), choose \(Z_0\) after the bounded-height
comparison, then fix the cone and barrier constants. Reduce \(q_1\)
so that
\begin{equation}
\label{eq:direct-contact-prefix-absorption-choice}
 q_1Z_0\le \frac{b_1}{4C_B},
 \qquad z_0=q_1^2Z_0\le \frac{q_1}{4C_B},
\end{equation}
where \(C_B\) bounds \(\lvert B_\eta\rvert\) on the comparison box and
\(b_1\) is the positive coefficient margin in
\eqref{eq:normal-sign-margins-new}. With \(z_*\) already fixed by the
common-cone construction, reduce \(\lvert q_2\rvert\) (and hence
\(d_{\rm c}\)) relative to this \(z_*\), after \(q_1\). The absorptions
below are then uniform in the independent limits of \(q_1\) and \(q_2\).

For \(\widehat q_0\le0\), choose initially \(\widehat h\le1/2\).
On \(h\le z\le z_0\), the terms \(B_\eta/|\eta_\sigma|\) in the
paired identity \eqref{eq:paired-log-derivative-new} are absorbed
by half of the positive terms \(q_1/(|\eta_\sigma|z)\): indeed,
\eqref{eq:direct-contact-prefix-absorption-choice} gives
\(z\le z_0\le q_1/(4C_B)\), so
\( |B_\eta|/|\eta_\sigma|\le q_1/(4|\eta_\sigma|z)\).
The fixed positive-radius bound
gives a gain at least \(c_Lt_h\). By
\eqref{eq:direct-contact-cap}, the cubic loss in this prefix is at most
\(C_Ld_{\rm c} t_h\). On \(z\ge z_0\), the lower cones bound the
\(B_\eta/|\eta_\sigma|\) integrals; the positive upper cone bounds
the cubic integral on the positive branch, and
\eqref{eq:direct-contact-tail-loss} bounds it on the negative branch.
Therefore
\[
 \log\frac{I_{-,h}}{I_{+,h}}
 \ge (c_L-C_Ld_{\rm c})t_h-C\sqrt{z_*}-Cd_{\rm c}\log(1+v_0)
 \ge \tfrac12c_Lt_h-C_L .
\]
The last step uses \eqref{eq:direct-contact-source-growth} and
\(\log(1+t_h)\le t_h\), reducing \(d_{\rm c}\) after \(L,Z_0\) and
\eqref{eq:direct-contact-prefix-absorption-choice} so that the
coefficient of \(t_h\) remains at least \(c_L/2\).
Choose a fixed \(\widehat h_{\rm lo}>0\)
so small that this bound exceeds \(\log(1/\Lambda_-)\) for
\(\widehat h\le \widehat h_{\rm lo}\). Then \(\Delta_h>0\) there.

For \(\widehat q_0\ge0\), the lower end of the passing interval is
controlled by the radial flux. Consider central heights with small positive flux
\(0<f\le1\). Then
\[
 \widehat h\le (L+1)/b_0,\qquad \widehat h\ge f/\sup B .
\]
For \(0\le t\le1\), subtracting the value of \(ZB\)
at the central point in \eqref{eq:direct-contact-equations} gives
\[
 (u_-^2/2)_t\le C_Lu_-+d_{\rm c}u_-^2+f+C_Lt .
\]
The explicit upper comparison
\[
 \bar v=K(t+\sqrt{ft}),\qquad
 (\bar v^2/2)'=K^2(t+\tfrac32\sqrt{ft}+f/2)
\]
proves \(u_-\le\bar v\). Choose \(K\) large first to dominate the
linear, \(f\), and \(t\) terms, then \(d_{\rm c}\) small to absorb
\(d_{\rm c}\bar v^2\). A fixed provisional radius box larger than
\(4K\) completes the first-contact continuation argument.
Thus the favorable negative-branch term in the paired identity
satisfies
\[
 \int_0^1\frac{dt}{u_-(t)}
 \ge \frac2K\log(1+f^{-1/2})
 \ge c_L\log(1/f)-C_L .
\]
Choose \(Z_0>\e(L+1)/b_0\), so this interval lies in the prefix.
The other prefix losses and the whole tail are bounded
as above. Since \(t_h\le\log(C_L/f)\), the result is
\begin{equation}
\label{eq:direct-contact-flux-gain}
 \log\frac{I_{-,h}}{I_{+,h}}
 \ge \tfrac12c_L\log(1/f)-C_L .
\end{equation}
Choose \(f_{\rm lo}>0\) so small that the right side exceeds
\(\log(1/\Lambda_-)\) for \(f\le f_{\rm lo}\). Thus \(\Delta_h>0\)
on this interval, whose connectedness follows from \(f_{\widehat h}>0\).
This estimate is uniform as \(f/\widehat h\to0\).

At the other end let \(\theta_0\) be the constant of
\cref{lem:relative-height-paired-new}, and choose
\(\widehat h_{\rm hi}>(L+1)/b_0+\widehat h_{\rm lo}\) so large that
\(L/\widehat h_{\rm hi}+\widehat h_{\rm hi}^{-1/2}\le\theta_0/2\).
For \(h\ge q_1^2\widehat h_{\rm hi}\),
\[
 |q_0|/h\le L/\widehat h_{\rm hi},\qquad
 q_1/\sqrt h\le \widehat h_{\rm hi}^{-1/2},
\]
and \(|q_2|\le\theta_0/2\) for small parameters. Hence
\cref{lem:relative-height-paired-new} applies and gives
\(\Delta_h<0\) at these central heights.

Write \(\widehat h_{\rm a}=\widehat h_{\rm lo}\) when \(\widehat q_0\le0\), and otherwise
let \(\widehat h_{\rm a}\) be the unique solution of \(f(\widehat h_{\rm a})=f_{\rm lo}\).
Put \(\widehat h_{\rm b}=\widehat h_{\rm hi}\). Taking \(f_{\rm lo}\le1\), the choice
of \(\widehat h_{\rm b}\) gives \(f(\widehat h_{\rm b})>1\) for \(\widehat q_0>0\), so \(\widehat h_{\rm a}<\widehat h_{\rm b}\) in both
cases. At central heights in \(\cH_p\), the displacement has
\(\Delta_h>0\) for \(\widehat h\le \widehat h_{\rm a}\) and \(\Delta_h<0\) for
\(\widehat h\ge \widehat h_{\rm b}\), including the height thresholds whenever they
are admissible.
The middle interval is \([\widehat h_{\rm a},\widehat h_{\rm b}]\). On bounded heights,
\[
 f(\widehat h)-(B_0\widehat h-\widehat q_0)
       =O(q_1^2\widehat h^2),\qquad B_0=B(0,0,p).
\]
Consequently the central heights of the middle interval belong, after
parameter reduction, to
\begin{equation}
\label{eq:direct-contact-compact-passing}
 \mathcal K_{\rm c}=\left\{(\widehat h,\widehat q_0,B_0):
 \begin{array}{c}
 b_0\le B_0\le B_{\max},\quad |\widehat q_0|\le L,\\
 h_*\le\widehat h\le\widehat h_{\rm hi},\quad
 B_0\widehat h-\widehat q_0\ge f_*
 \end{array}\right\},
\end{equation}
where
\[
 h_*:=\min\{\widehat h_{\rm lo},f_{\rm lo}/(2B_{\max})\}>0,
 \qquad f_*:=\min\{b_0\widehat h_{\rm lo},f_{\rm lo}/2\}>0.
\]
Here \(B_{\max}\) is a fixed upper bound for \(B_0\).
Every fixed \((\widehat q_0,B_0)\) fibre of \(\mathcal K_{\rm c}\) is an interval,
possibly empty. This compact set lies strictly inside the model passing domain
\begin{equation}
\label{eq:direct-contact-model-domain}
 |\widehat q_0|\le L,\qquad b_0\le B_0\le B_{\max},\qquad
 \widehat h>\max\{0,\widehat q_0/B_0\}.
\end{equation}
For the estimates below, choose one fixed
\(\widehat h_{\rm ref}>\widehat h_{\rm hi}+1+2L/b_0\),
and let $\mathcal K_{\rm c}^+$ be the compact set obtained from $\mathcal K_{\rm c}$
by replacing its upper height bound by $\widehat h_{\rm ref}+1$
and its lower height and flux bounds by $h_*/2$ and $f_*/2$.
Every fibre contains the reference height $\widehat h_{\rm ref}$ and
is still an interval. By the estimate of \(f\) above, the radial flux
at the central point is positive on a neighborhood of this enlarged
set for small parameters.

\subsubsection{Scattering and a strictly decreasing derivative ratio}

The outer signs leave only the compact passing set
\(\mathcal K_{\rm c}\). On this set we compare the derivatives of the
limiting scattering maps, and then transfer their strict ratio
margin to the endpoint maps.
For fixed contact heights the limiting equations are
\begin{equation}
\label{eq:direct-contact-model}
 (\mathcal U_\sigma^0)_Z=B_0-\frac{\widehat q_0+\sigma\sqrt{2\mathcal U_\sigma^0}}{Z},
 \qquad \mathcal U_\sigma^0(\widehat h;\widehat h)=0 .
\end{equation}
Define
\begin{equation}
\label{eq:direct-contact-scattering}
 S_\sigma(\widehat h;\widehat q_0,B_0)=
 \lim_{Z\to\infty}
 [\mathcal U_\sigma^0(Z;\widehat h_{\rm ref})-\mathcal U_\sigma^0(Z;\widehat h)] .
\end{equation}
The limit exists with local $C^2$ convergence of order $O(Z^{-1/2})$
in $\widehat h$, uniformly on compact passing sets. Fix a height beyond
all the rescaled central heights of such a set. The time equations
\eqref{eq:direct-contact-time} with
$q_1=q_2=0$ have positive initial radial flux, and their solutions
meet this height section transversely. The resulting
finite-height maps are analytic.
The limiting barriers put both graphs in $cZ\le \mathcal U_\sigma^0\le CZ$
above a sufficiently large fixed cut. There
\[
 \partial_{\mathcal U}\mathcal G_\sigma^0=-\frac{\sigma}{Z\sqrt{2\mathcal U_\sigma^0}},\qquad
 |\partial_{\mathcal U}\mathcal G_\sigma^0|\le CZ^{-3/2},\qquad
 |\partial_{\mathcal U}^2\mathcal G_\sigma^0|\le CZ^{-5/2}.
\]
The mean-value equation for two graphs and the triangular first- and
second-variation equations have these integrable coefficients. Hence
the graph difference and its first two \(\widehat h\)-derivatives remain
bounded, and their \(Z\)-derivatives are \(O(Z^{-3/2})\). Integration to
infinity gives the asserted \(C^2\) convergence and rate. The same
majorants give continuous dependence on \((\widehat q_0,B_0)\), and
\(S_\sigma(\widehat h_{\rm ref})=0\).

The variation formula for a moving initial point gives
\begin{equation}
\label{eq:direct-contact-model-derivative}
 S_\sigma'(\widehat h)=
 \left(B_0-\frac{\widehat q_0}{\widehat h}\right)
 \exp\left[-\sigma\int_{\widehat h}^\infty
          \frac{dZ}{Z\sqrt{2\mathcal U_\sigma^0(Z;\widehat h)}}\right]>0 .
\end{equation}
Indeed, the initial factor of $\partial_{\widehat h}\mathcal U_\sigma^0$
is $-(B_0-\widehat q_0/\widehat h)$ for both signs. The integral in
\eqref{eq:direct-contact-model-derivative} converges both at the
central point and at infinity.

\begin{samepage}
\begin{lemma}[Scattering derivative ratio]
\label{lem:direct-contact-scattering-ratio}
For fixed $(\widehat q_0,B_0)$, the function
\[
 \mathscr R(\widehat h):=S_-'(\widehat h)/S_+'(\widehat h)
\]
is strictly decreasing throughout the passing interval
\eqref{eq:direct-contact-model-domain}. On each fixed compact
passing set there are constants $c_1,c_2>0$ such that
\begin{equation}
\label{eq:direct-contact-ratio-margins}
 S_+'\ge c_1,\qquad -\mathscr R'\ge c_2.
\end{equation}
\end{lemma}
\end{samepage}

\begin{proof}
Put $t=\log(Z/\widehat h)$ and
$W_\sigma(t;\widehat h)=\mathcal U_\sigma^0(\widehat h \e^t;\widehat h)$.
The initial time is now fixed, and
\[
 (W_\sigma)_t=B_0\widehat h \e^t-\widehat q_0-\sigma\sqrt{2W_\sigma},
 \qquad W_\sigma(0;\widehat h)=0.
\]
Take $V_\sigma=\partial_{\widehat h}W_\sigma$ at fixed $t$.
With $f^0=B_0\widehat h-\widehat q_0>0$,
the regular time equations at the central point give, uniformly on compact
passing sets,
\[
 W_\sigma=f^0t+O(t^{3/2}),\qquad
 V_\sigma=B_0t+O(t^{3/2}).
\]
For example, in the analytic time variable $\tau$,
$u_\sigma=f^0\tau+O(\tau^2)$ and
$t=f^0\tau^2/2+O(\tau^3)$; inversion in $\sqrt t$ gives these
expansions with their derivatives in $\widehat h$.
For $t>0$ the variation equation is
\[
 (V_\sigma)_t=B_0\e^t-\frac{\sigma V_\sigma}{\sqrt{2W_\sigma}},
 \qquad V_\sigma(0)=0.
\]
Its coefficient is $O(t^{-1/2})$ at zero. Variation of constants,
first from positive time and then by passage to zero, yields
\begin{equation}
\label{eq:direct-contact-fixed-time-variation}
 V_\sigma(t)=B_0\int_0^t \e^{\bar t}
 \exp\left[-\sigma\int_{\bar t}^t\frac{du}{\sqrt{2W_\sigma(u)}}\right]d\bar t>0.
\end{equation}
For large $t$, the cone estimates give $W_\sigma\asymp \e^t$.
The tail coefficient of this variation equation is integrable,
so \eqref{eq:direct-contact-fixed-time-variation} gives
$V_\sigma=O(\e^t)$. Thus both
$(2W_\sigma)^{-1/2}$ and
$V_\sigma(2W_\sigma)^{-3/2}$ are bounded by
$Ct^{-1/2}$ near zero and by $C\e^{-t/2}$ at infinity,
uniformly on every compact passing set. Dominated differentiation gives
\[
 \mathcal T_\sigma(\widehat h):=\int_0^\infty
            \frac{dt}{\sqrt{2W_\sigma(t;\widehat h)}},\qquad
 \mathcal T_\sigma'=-\int_0^\infty
       \frac{V_\sigma(t;\widehat h)}{(2W_\sigma(t;\widehat h))^{3/2}}\,dt<0,
\]
and $\mathcal T_\sigma'$ is continuous in the passing parameters.
The common positive factor $B_0-\widehat q_0/\widehat h$ in
\eqref{eq:direct-contact-model-derivative} cancels, giving
\[
 \mathscr R=\exp(\mathcal T_-+\mathcal T_+),\qquad
 \mathscr R'=\mathscr R(\mathcal T_-'+\mathcal T_+')<0.
\]
Continuity and compactness give \eqref{eq:direct-contact-ratio-margins}.
\end{proof}

\subsubsection{Endpoint preparation and the two-zero bound}

To transfer the scattering comparison to the endpoint maps, we must
retain the quadratic tail factor even when its logarithm is unbounded.
The common endpoint amplitude is
\begin{equation}
\label{eq:direct-contact-amplitude}
 E=q_1^2(z_*/q_1^2)^{-2q_2}.
\end{equation}
By \eqref{eq:direct-tail-amplitude} with \(t=q_1^2\),
\(E\le z_*^{1/4}q_1^{3/2}\to0\) for \(|q_2|<1/8\), while
\(q_2\log q_1\) may be unbounded.
Let \(b_\sigma=I_\sigma(q_1^2\widehat h_{\rm ref})\) and set
\[
 A_\sigma(p)=
 \frac{\partial_{\widehat h} I_\sigma(q_1^2\widehat h)|_{\widehat h=\widehat h_{\rm ref}}}
      {E S_\sigma'(\widehat h_{\rm ref})}.
\]
On the enlarged compact passing set above,
\begin{equation}
\label{eq:direct-contact-actual-preparation}
 I_\sigma(q_1^2\widehat h)=b_\sigma+
 E A_\sigma(p)[S_\sigma(\widehat h)+R_\sigma(\widehat h,p)],\qquad
 \|R_\sigma\|_{C^2}\longrightarrow0,
\end{equation}
where \(A_\sigma\) and \(A_\sigma^{-1}\) are bounded and positive.
The signs \(\sigma=-\) and \(\sigma=+\) label the negative and positive
branches, as in \eqref{eq:direct-contact-equations}. The coefficient
ratio satisfies
\begin{equation}
\label{eq:direct-contact-slope-quotient}
 A_-/A_+\le \exp(2C_{\rm cone}\sqrt{z_*})+o(1).
\end{equation}
Choose a large fixed matching height \(Z_c>2\widehat h_{\rm ref}\), and put
\[
 D_{\sigma,c}(\widehat h)
 =\mathcal U_\sigma(Z_c;\widehat h_{\rm ref})-\mathcal U_\sigma(Z_c;\widehat h).
\]
At each fixed \(Z_c\), the finite-height analytic flow
\eqref{eq:direct-contact-time} converges, with all state derivatives
through any prescribed finite order, to \eqref{eq:direct-contact-model}. The
scattering estimate therefore gives
\[
 \|D_{\sigma,c}-S_\sigma\|_{C^2}
 \le CZ_c^{-1/2}+C(Z_c)(q_1+|q_2|).
\]
After fixing \(Z_c\) and reducing the parameters, these differences
have bounded $C^2$ norms. Integration of the comparison equation gives
\(\mathcal U_\sigma^0(Z;\widehat h)=B_0Z+O(\sqrt Z+\log Z)\), uniformly on
each fixed compact set of rescaled central heights. Since
\(B_0\in[b_0,\sup B]\) lies in the interior of the common cone, a
sufficiently large \(Z_c\) puts the normalized energies
\(\mathcal U_\sigma^0(Z_c;\widehat h)/Z_c\), and after parameter reduction also
\(\mathcal U_\sigma(Z_c;\widehat h)/Z_c\), in that cone. At \(t=q_1^2Z_c\),
moreover,
\[
 |q_0|/t\le L/Z_c,\qquad q_1/\sqrt t=Z_c^{-1/2}.
\]
Apply \cref{lem:direct-weighted-tail} at \(t=q_1^2Z_c\), retaining
\((z_*/t)^{-2q_2}\). Its centered map is
\(a_{\sigma,c}x+N_\sigma(x)\), where
\(N_\sigma(0)=N_\sigma'(0)=0\),
\(\|N_\sigma\|_{C^4}\le Ce(t)\), and
\(e(t)=q_1\sqrt{Z_c}+Z_c^{-1/2}\). Here
\(x=-D_{\sigma,c}/Z_c=O_{C^2}(Z_c^{-1})\), so Taylor's formula gives
\(\|N_\sigma(x)\|_{C^2}\le Ce(t)/Z_c^2\). Since
\(t(z_*/t)^{-2q_2}/E=Z_c^{1+2q_2}\), the resulting endpoint is
\[
 \frac{I_\sigma(q_1^2\widehat h)-b_\sigma}{E}
 =Z_c^{2q_2}a_{\sigma,c}
       [D_{\sigma,c}(\widehat h)+n_{\sigma,c}(\widehat h)],\qquad
 \|n_{\sigma,c}\|_{C^2}\le Ce(t)/Z_c,
\]
with \(n_{\sigma,c}=n_{\sigma,c}'=0\) at \(\widehat h_{\rm ref}\).
Normalize by the positive reference derivative
\[
 A_\sigma=Z_c^{2q_2}a_{\sigma,c}
 \frac{D_{\sigma,c}'(\widehat h_{\rm ref})}
      {S_\sigma'(\widehat h_{\rm ref})}.
\]
The scattering margin bounds the denominator below, and the preparation
\eqref{eq:direct-contact-actual-preparation} follows with
\[
 \|R_\sigma\|_{C^2}
 \le CZ_c^{-1/2}+C(Z_c)(q_1+|q_2|)
           +C(q_1/\sqrt{Z_c}+Z_c^{-3/2}).
\]
Choose \(Z_c\) large and then \(q_1,|q_2|\) small. With this fixed
\(Z_c\), \(q_2\log Z_c\to0\), so \(A_\sigma^{\pm1}\) are bounded.

Since the normalized energies at \(t\) lie in the common cone,
\eqref{eq:direct-tail-first-slope} gives
\[
 |\log a_{\sigma,c}|\le C_{\rm cone}\bigl(\sqrt{z_*}+Z_c^{-1/2}\bigr).
\]
The factor \(Z_c^{2q_2}\) cancels from \(A_-/A_+\), and
the two reference derivative ratios tend to one in the
successive choices above, proving \eqref{eq:direct-contact-slope-quotient}.

We now combine this preparation with the outer derivative signs on
the same admissible height interval. Let
\(\widehat{\mathcal H}_p=q_1^{-2}\cH_p\) be its rescaling in the
coordinate \(\widehat h\). If \(\widehat{\mathcal H}_p\cap[\widehat h_{\rm a},\widehat h_{\rm b}]\) is empty, connectedness puts
\(\widehat{\mathcal H}_p\) in one outer interval, where the strict derivative sign already
gives at most one zero. Otherwise, the endpoint graphs are defined on
the whole fibre of $\mathcal K_{\rm c}^+$, and every endpoint on this fibre, including
the reference endpoint, is within $O(E)$ of the corresponding
admissible endpoint. By \cref{lem:exterior-multiplier-new}, all these
endpoints lie in \(\widetilde W_\pm\) for small parameters.
By \eqref{eq:direct-contact-slope-quotient} and the choice of \(z_*\),
a reduction of parameters gives
\[
 0<\lambda_0\le\lambda(p):=
 \Pi_p'(b_-)A_-/A_+\le\bar\lambda<1.
\]
Put $G(\widehat h;\widehat q_0,B_0,\lambda)=\lambda S_-(\widehat h)-S_+(\widehat h)$.
Taylor's formula for the exterior map and
\eqref{eq:direct-contact-actual-preparation} give
\begin{equation}
\label{eq:direct-contact-completed-preparation}
 \frac{\Delta(q_1^2\widehat h)}{EA_+}
 =D(p)+G(\widehat h;\widehat q_0,B_0,\lambda(p))+R_\Delta(\widehat h,p),
 \qquad \|R_\Delta\|_{C^2}\to0 .
\end{equation}
Here \(D=(\Pi_p(b_-)-b_+)/(EA_+)\).
The exterior Taylor remainder is
\(O_{C^2}(E)\), since the normalized endpoint functions have
bounded \(C^2\) norms and the exterior map has bounded
derivatives through order four.

Write \(\widehat\Delta=\Delta(q_1^2\widehat h)/(EA_+)\).
Apply \cref{lem:direct-derivative-ratio-count} on each fibre of
\(\mathcal K_{\rm c}\) to \(-\widehat\Delta\), with
\(\phi=S_+\), \(\psi=S_-\) and
\(\lambda\in[\lambda_0,\bar\lambda]\). The uniform margins of
\cref{lem:direct-contact-scattering-ratio} give
\begin{equation}
\label{eq:direct-contact-completed-ratio}
 \left(\frac{\widehat\Delta'}{S_+'}\right)'<0 .
\end{equation}
In increasing \(T\), the large-height, compact and small-height
intervals are ordered by the cuts
\(\widehat h_{\rm b}\) and \(\widehat h_{\rm a}\). The outer signs
exclude slope zeros there, including at the cuts. In the middle,
\eqref{eq:direct-completed-log-ratio} writes
\(\widehat\Delta'/S_+'\) as \(\e^{\mathcal L}-1\) times the positive
factor \(q_1^2I_{+,h}/(EA_+S_+')\). Thus the displayed inequality gives
\(\mathcal L_{\widehat h}<0\), or \(\mathcal L_T>0\), at every slope
zero. All preparations represent the same displacement on the enlarged
windows. The second case of \cref{cor:ordered-cover-zero-count} gives
at most two zeros on the whole interval.

This proves \cref{prop:direct-bounded-contact}, with
\[
 N_{\rm ct}(L)=2.
\]
